\begin{tikzpicture}[
  x=1mm,y=1mm,
  >=Stealth,
  font=\figurefont\fontsize{8}{10}\selectfont
]

\tikzset{
  vA/.style={circle,draw=FigureBlue,fill=FigureBlueFill,minimum size=4.6mm,inner sep=0pt,line width=.55pt,text=FigureInk},
  vT/.style={circle,draw=FigureYellow,fill=FigureYellowFill,minimum size=4.6mm,inner sep=0pt,line width=.6pt,text=FigureInk},
  vB/.style={circle,draw=FigureGreen,fill=FigureGreenFill,minimum size=4.6mm,inner sep=0pt,line width=.55pt,text=FigureInk},
  edge/.style={draw=FigureStroke,line width=.55pt},
  activeedge/.style={draw=FigureRed,line width=.8pt},
  note/.style={font=\FigureNoteFont\fontsize{7.5}{9}\selectfont,text=FigureInk},
  muted/.style={font=\FigureNoteFont\fontsize{7.0}{8.5}\selectfont,text=FigureMuted}
}

% (a) Separator T and 3-value trial vector
\begin{scope}[xshift=0mm]
  \draw[rounded corners=2mm,dashed,draw=FigureBlue,fill=FigureBlueFill!35] (-3,-1) rectangle (13,19);
  \draw[rounded corners=2mm,dashed,draw=FigureYellow,fill=FigureYellowFill!45] (19,1) rectangle (29,17);
  \draw[rounded corners=2mm,dashed,draw=FigureGreen,fill=FigureGreenFill!35] (35,-1) rectangle (51,19);

  \node[vA] (a1) at (2,14) {};
  \node[vA] (a2) at (8,9) {};
  \node[vA] (a3) at (2,4) {};

  \node[vT] (t1) at (24,13) {};
  \node[vT] (t2) at (24,5) {};

  \node[vB] (b1) at (39,13) {};
  \node[vB] (b2) at (46,9) {};
  \node[vB] (b3) at (39,5) {};

  % Internal edges (diff = 0)
  \draw[edge] (a1)--(a2)--(a3)--(a1);
  \draw[edge] (b1)--(b2)--(b3);
  \draw[edge] (t1)--(t2);

  % Cross edges to T (active energy)
  \draw[activeedge] (a1)--(t1);
  \draw[activeedge] (a2)--(t1);
  \draw[activeedge] (a2)--(t2);
  \draw[activeedge] (a3)--(t2);
  \draw[activeedge] (t1)--(b1);
  \draw[activeedge] (t2)--(b3);

  \node[note,text=FigureBlue] at (5,21.5) {$A$：$x_i=+|B|$};
  \node[note,text=FigureYellow] at (24,21.5) {$T$：$x_i=0$};
  \node[note,text=FigureGreen] at (43,21.5) {$B$：$x_i=-|A|$};
  \node[muted,text=FigureRed] at (24,-3.5) {仅连向分隔集 $T$ 的红边产生能量};
  \node[note] at (24,-8.5) {(a) 点分隔集 $T$（$|T|=\kappa$）上的三值试探向量};
\end{scope}

% (b) P_4 vs C_4 verification
\begin{scope}[xshift=60mm]
  % Top: P_4
  \node[vA] (p1) at (0,14) {$1$};
  \node[vT] (p2) at (11,14) {$2$};
  \node[vB] (p3) at (22,14) {$3$};
  \node[vB] (p4) at (33,14) {$4$};
  \draw[activeedge] (p1)--(p2);
  \draw[activeedge] (p2)--(p3);
  \draw[edge] (p3)--(p4);
  \node[muted] at (16.5,19) {$P_4$：删点 $2$ 即断（$\kappa=\lambda=\delta=1$）};
  \node[muted] at (16.5,9.5) {$\lambda_2=2-\sqrt 2\approx 0.59 < 1$};

  % Bottom: C_4
  \node[vA] (c1) at (6,2) {$1$};
  \node[vT] (c2) at (16.5,5.5) {$2$};
  \node[vB] (c3) at (27,2) {$3$};
  \node[vT] (c4) at (16.5,-1.5) {$4$};
  \draw[activeedge] (c1)--(c2)--(c3)--(c4)--(c1);
  \node[muted] at (16.5,-5) {$C_4$：删 $\{2,4\}$ 断开，$\lambda_2=\kappa=\lambda=\delta=2$};
  \node[note] at (16.5,-9.5) {(b) $P_4$ 与 $C_4$ 的连通度对比};
\end{scope}
\end{tikzpicture}
